% Generated draft: requires mathematical and source review.
\documentclass[11pt]{article}
\usepackage[margin=25mm]{geometry}
\usepackage{fontspec}
\setmainfont{Latin Modern Roman}
\usepackage{amsmath,amssymb,mathtools,unicode-math}
\setmathfont{Latin Modern Math}
\usepackage{enumitem,needspace,xurl,hyperref}
\setlength{\parindent}{0pt}
\setlength{\parskip}{5pt}
\setlength{\emergencystretch}{3em}
\newcommand{\N}{\mathbb{N}}
\newcommand{\Z}{\mathbb{Z}}
\newcommand{\Q}{\mathbb{Q}}
\newcommand{\R}{\mathbb{R}}
\newcommand{\C}{\mathbb{C}}
\newcommand{\F}{\mathbb{F}}
\newcommand{\A}{\mathbb{A}}
\newcommand{\PP}{\mathbb{P}}
\newcommand{\E}{\mathbb{E}}
\newcommand{\eps}{\varepsilon}
\newcommand{\dd}{\,\mathrm{d}}
\begin{document}
\section*{UnsolvedMath ID: 7500119}
\textbf{Source problem code:} catalogue entry
\subsection{Definitions and mathematical statement}
The base theory \(\mathsf{RCA}_0\) is the standard subsystem of second-order arithmetic with elementary arithmetic axioms, \(\Sigma^0_1\)-induction and \(\Delta^0_1\)-comprehension, allowing set parameters. The system \(\Pi^1_1\text{-}\mathsf{CA}_0\) adds comprehension for \(\Pi^1_1\) formulas: for every such \(\varphi(n,\vec X,\vec a)\),
\[
\forall\vec X\,\forall\vec a\ \exists A\ \forall n,
\qquad n\in A\iff\varphi(n,\vec X,\vec a).
\]
Here a \(\Pi^1_1\) formula has universal set quantifiers followed by an arithmetical formula, with number quantifiers and parameters allowed.

Write \(\mathsf{FZ}\) for the intended second-order-arithmetic formulation of the Furstenberg--Zimmer structure theorem. The exact reverse-mathematical question is
\[
\mathsf{RCA}_0+\mathsf{FZ}\ \vdash\
\Pi^1_1\text{-comprehension},
\]
meaning that every instance of the displayed comprehension scheme is provable. This is the full implication over the stated base, not merely an assertion about standard or omega-models and not an equivalence claim.

Classically, the Furstenberg--Zimmer structure theorem describes measure-preserving systems/extensions using relatively compact or distal structure, transfinite compact-extension towers, and a relatively weak-mixing remainder. A probability-preserving system can be presented as \((X,\mathcal B,\mu,T)\) with invertible measure-preserving \(T\), but a reverse-mathematical sentence requires precise countable codes for the systems, factors, measures, and any tower/well-order data. The record does not specify those codes or the theorem's exact version. They cannot be replaced by an arbitrary classical formulation when the proof-theoretic strength is the object of the question.

The supplied research summary reports no verified implication or resolution in the accessible literature as of mid-2026. No upper-bound theorem or lower-bound reversal is supplied.
\subsection{Short English statement}
Does the intended formal Furstenberg--Zimmer structure theorem, added to \(\mathsf{RCA}_0\), prove every \(\Pi^1_1\)-comprehension instance? The precise countable coding of the structure theorem must be fixed.
\subsection{Sources}
\begin{itemize}
\item Local UnsolvedMath registry, permanent ID \(7500119\), exact source code \texttt{catalogue entry}, ``Furstenberg--Zimmer Structure Theorem.'' Proposed year 2011; no proposer is supplied.
\item Montalbán, \emph{Open Questions in Reverse Mathematics}, Question 19; \url{https://math.berkeley.edu/~antonio/papers/questionsRM.pdf}. Recorded access 2026-07-29; author-preprint text with manual normalization.
\item The 2026-07-29 nested-list audit records appears\_open with formulation still NEEDS\_REVIEW. The supplied OPEN-TRIAGE research summary reports no verified solution as of mid-2026, without an additional precise publication citation.
\end{itemize}
\end{document}
